Define \(\lambda_k=c_0/r\).  Because the intervals \(I_1,\ldots,I_r\) partition \([0,1]\) up to endpoints of Lebesgue measure zero,
\[
 \sum_{k=1}^r\lambda_k\psi_k(u)\psi_k(v)
 =c_0\mathbf1\{u\text{ and }v\text{ belong to the same }I_k\}
 =W_0(u,v)
\tag{1}
\]
for almost every \((u,v)\).  Moreover,
\[
 \int_0^1\psi_k(u)\psi_\ell(u)\,du=\mathbf1\{k=\ell\}.
\tag{2}
\]
All \(r\) eigenvalues in (1) are nonzero, and the definition of \(c_0\) gives
\[
 |\lambda_k|=c_0/r\geq\underline\lambda,\qquad
 \int_0^1W_0(u,v)\,dv=c_0/r\geq\kappa^{-1}.
\tag{3}
\]
The compatibility condition \(c_0\leq\kappa\) gives \(0\leq W_0\leq\kappa\).  Together with \({\rm ass:edge\text{-}probability}\), this makes every edge probability \(\rho_nW_0(U_i,U_j)\) lie in \([0,1]\).

To verify the Bernstein member property, note that
\[
 \overline\psi_k:=\int_0^1\psi_k(u)\,du=1/\sqrt r.
\]
For \(r\geq2\), the centered variable \(X_k=\psi_k(U)-\overline\psi_k\) takes the values \((r-1)/\sqrt r\) and \(-1/\sqrt r\), and hence \(|X_k|\leq(r-1)/\sqrt r\leq b_\psi\).  For every integer \(m\geq2\),
\[
 \mathbb E|X_k|^m
 \leq b_\psi^{m-2}\mathbb EX_k^2
 \leq\frac{m!}{2}b_\psi^{m-2}\mathbb EX_k^2.
\tag{4}
\]
For \(r=1\), \(X_k=0\), so (4) is immediate.  Thus the declared Bernstein moment condition holds.  Equations (1)--(4), the independent uniform latent types, \({\rm ass:graphon\text{-}draw}\), and \({\rm ass:bernoulli\text{-}design}\) establish the advertised known-rank equal-block graph submodel, including rank, positivity, degree regularity, spectral strength, and treatment overlap.

Put \(s_0=B/4\).  The function
\[
 p_{s_0}(t)=s_0^{-1}\cos^2\!\left(\frac{\pi t}{2s_0}\right)
 \mathbf1\{|t|\leq s_0\}
\tag{5}
\]
is nonnegative, is supported on \([-B/4,B/4]\), and integrates to one because the change of variables \(y=t/s_0\) gives
\[
 \int_{-s_0}^{s_0}p_{s_0}(t)\,dt
 =\int_{-1}^{1}\cos^2(\pi y/2)\,dy=1.
\tag{6}
\]
Let \(\xi_i\) be independent draws from this density, independent of the graph primitives.  For a deterministic function \(g\), define
\[
 f_i^\sigma(a,x)=\xi_i+\sigma g(x),\qquad \sigma\in\{-1,1\}.
\tag{7}
\]
The resulting \((U_i,f_i^\sigma)\) are independent and identically distributed.  Since \(g\) is independent of \(a\), \({\rm def:target}\) gives
\[
 \theta_n^\sigma=\sigma g'(1/2).
\tag{8}
\]

We now specialize the explicit response experiments proved in \({\rm thm:unrestricted\text{-}supported\text{-}converse}\) to the graphon \(W_0\).  Write
\[
 \phi(t)=t(1-t^2)^4\mathbf1\{|t|\leq1\}.
\tag{9}
\]
Its zero extension is \(C^3\), because the polynomial has a zero of order four at each endpoint, and \(\phi'(0)=1\).  For \(K_j=\sup_t|\phi^{(j)}(t)|\), choose fixed \(\gamma>0\) such that \(\gamma K_3\leq1\) and sufficiently small for the testing bounds below.  The three functions are
\[
\begin{aligned}
 g_G(x)&=a_G(x-1/2),
 &a_G&=\zeta s_0\sqrt{d/n},\\
 g_I(x)&=\gamma Lh_0^3\phi((x-1/2)/h_0),&&\\
 g_L(x)&=\gamma Lh_L^3\phi((x-1/2)/h_L),
 &h_L&=(4Kd)^{-1},
\end{aligned}
\tag{10}
\]
where \(\zeta>0\) is a sufficiently small fixed constant and \(K>\kappa\) is fixed and chosen so that
\[
 C_{\log}\{K\log(K/\kappa)-K+\kappa\}\geq4.
\tag{11}
\]
The interior function \(g_I\) is used only when \(L>0\) and \(d^{-1}<h_0<d^{-1/2}\); when \(L=0\), only the graph pair is needed.

We verify that (7) stays in the same response class.  For \(g_G\), \(d/n\to0\) implies eventually
\[
 \|g_G\|_\infty\leq B-s_0,\qquad |g_G'|=a_G\leq B,\qquad
 g_G''=g_G'''=0.
\tag{12}
\]
For either local bump \(g(x)=\gamma Lh^3\phi((x-1/2)/h)\), differentiation gives
\[
 \|g^{(j)}\|_\infty\leq\gamma Lh^{3-j}K_j,\qquad 0\leq j\leq3.
\tag{13}
\]
In the interior regime, \(h_0<d^{-1/2}\to0\); for the lattice bump, \(h_L\to0\).  Hence (13), \(\gamma K_3\leq1\), and fixed \(B,L\) imply, for all large \(n\),
\[
 \|g\|_\infty\leq B-s_0,\qquad
 \max_{1\leq j\leq2}\|g^{(j)}\|_\infty\leq B,\qquad
 \|g'''\|_\infty\leq L.
\tag{14}
\]
Equations (5), (7), (12), and (14) verify the response envelope and the third-derivative radius.  Thus every pair in (10), with the retained \(W_0\) graph law, belongs to \(\mathcal M_{n,d,L}\) for all sufficiently large \(n\).

For clarity, we recall the quantitative two-point conclusion established in \({\rm thm:unrestricted\text{-}supported\text{-}converse}\).  If \(P_g^+\) and \(P_g^-\) denote the two full observed-data laws generated by (7), its Hellinger calculation gives
\[
 H^2(P_g^+,P_g^-)
 \leq\frac{\pi^2}{s_0^2}\sum_{i=1}^n\mathbb E[g(T_i)^2],
\tag{15}
\]
where \(T_i=M_i/N_i\) for \(N_i>0\) and \(T_i=0\) for isolates.  Put \(\eta=(1-2\alpha)/2\in(0,1)\).  Whenever the right side of (15) is at most \(\eta^2\), total variation is at most \(\eta\), and the testing inequalities proved there yield, with \(a=|g'(1/2)|\),
\[
 \inf_{\widetilde\theta}\max_{\sigma\in\{-1,1\}}
 \mathbb E_{P_g^\sigma}|\widetilde\theta-\theta_n^\sigma|
 \geq\frac a2(1-\eta),
\tag{16}
\]
and, for every interval having coverage at least \(1-\alpha\) under both laws,
\[
 \max_{\sigma\in\{-1,1\}}
 \mathbb E_{P_g^\sigma}[\operatorname{length}(\mathcal C)]
 \geq a(1-2\alpha).
\tag{17}
\]
These are internal proved consequences of the declared converse theorem, not additional assumptions.

For the block graphon, conditional degrees have means bounded above and below by fixed multiples of \(d\).  The moment calculation in \({\rm thm:unrestricted\text{-}supported\text{-}converse}\) therefore gives
\[
 \mathbb E(T_i-1/2)^2\leq C_1/d.
\tag{18}
\]
Equations (10), (15), and (18) show
\[
 H^2(P_{g_G}^+,P_{g_G}^-)
 \leq C_2\zeta^2.
\tag{19}
\]
Choosing \(\zeta\) once so that (19) is at most \(\eta^2\), equations (8), (10), and (16)--(17) give lower bounds of order \(\sqrt{d/n}\) for both criteria.

The binomial small-ball calculation in the declared converse theorem gives, uniformly for \(d^{-1}\leq h\leq d^{-1/2}\),
\[
 \Pr\{|T_i-1/2|\leq h\}\leq C_3h\sqrt d.
\tag{20}
\]
Since \(nL^2h_0^7\sqrt d=1\), equations (10), (15), and (20) imply
\[
 H^2(P_{g_I}^+,P_{g_I}^-)
 \leq C_4\gamma^2,\qquad
 |g_I'(1/2)|=\gamma Lh_0^2
 =\gamma L^{3/7}(n^2d)^{-1/7}.
\tag{21}
\]
Choose \(\gamma\) small enough that the first quantity in (21) is at most \(\eta^2\).  Equations (16)--(17) then give the interior lower term for both criteria.

For the lattice bump, every nonzero centered exposure-lattice point of a unit with \(1\leq N_i\leq Kd\) has absolute value at least \((2N_i)^{-1}\geq(2Kd)^{-1}>h_L\), while \(g_L(1/2)=0\).  Thus \(g_L(T_i)=0\) for these units and also for isolates.  The binomial Chernoff calculation in the declared converse theorem and (11) give
\[
 \Pr(N_i>Kd)\leq n^{-4}.
\tag{22}
\]
Consequently,
\[
 H^2(P_{g_L}^+,P_{g_L}^-)
 \leq C_5nL^2d^{-6}n^{-4}=o(1),\qquad
 |g_L'(1/2)|=\gamma Lh_L^2
 =\frac{\gamma L}{16K^2d^2}.
\tag{23}
\]
For all sufficiently large \(n\), equations (16)--(17) give the degree-lattice lower term \(Ld^{-2}\).  This pair agrees exactly on every vertex of degree at most \(Kd\); only the uniformly negligible upper-degree tail contributes to (23).

It remains to compare the three separations with the analytic envelope.  Put \(c_{n,d}=n^{-1/2}d^{-1/4}\) and \(g_0(h)=Lh^2+c_{n,d}h^{-3/2}\).  For \(L>0\),
\[
 Lh_0^2=c_{n,d}h_0^{-3/2}.
\tag{24}
\]
At all large indices, \(\mathcal H_{n,d}=[d^{-1},d^{-1/2}]\).  Evaluation at the appropriate endpoint or at \(h_0\) gives
\[
 \inf_{h\in\mathcal H_{n,d}}g_0(h)
 \leq
 \begin{cases}
 2Ld^{-2},&h_0\leq d^{-1},\\
 2Lh_0^2,&d^{-1}<h_0<d^{-1/2},\\
 2\sqrt{d/n},&h_0\geq d^{-1/2}.
 \end{cases}
\tag{25}
\]
For \(L=0\), evaluation at \(d^{-1/2}\) gives the last order directly.  By \({\rm def:frontier\text{-}envelope}\), \(\mathfrak r^{\mathrm{sup}}_{n,d,L}=\sqrt{d/n}+\inf_{h\in\mathcal H_{n,d}}g_0(h)\).  The largest applicable separation in (19), (21), and (23) is therefore at least a fixed multiple of \(\mathfrak r^{\mathrm{sup}}_{n,d,L}\), because the maximum of two nonnegative quantities is at least half their sum.  Since all testing laws are in the explicit block submodel, equations (16)--(17) and \({\rm def:minimax\text{-}criteria}\) yield
\[
 R^*_{n,d,L}\geq c\mathfrak r^{\mathrm{sup}}_{n,d,L},
 \qquad
 \Lambda^*_{n,d,L}(\alpha)\geq c\mathfrak r^{\mathrm{sup}}_{n,d,L}
\tag{26}
\]
eventually, where \(c>0\) depends only on the fixed class constants and \(\alpha\), not on \(n,d\).  The compatibility inequalities were used only to certify the convenient block witness through (3)--(4).  Independently, \({\rm thm:unrestricted\text{-}supported\text{-}converse}\) proves (26) by retaining the graph primitives of any admissible member, so these sufficient block inequalities do not characterize nonemptiness of the maintained class.